\documentclass[11pt]{article}
\usepackage[margin=2.5cm]{geometry}
\usepackage{amsmath,amssymb,booktabs,array}

\begin{document}

\section*{Zone specification for peninsular electricity demand}

We regress hourly log peninsular electricity demand on daily heating and
cooling degree-days of four climate zones, with a separate coefficient for
each zone and hour of the day:
\begin{equation}
\begin{split}
\ln D_t = {} & \sum_{z \in \mathcal{Z}} \sum_{k=0}^{23}
  \beta^{C}_{z,k}\, \mathbf{1}\{h(t)=k\}\, CDD_{z,d(t)}
  + \sum_{z \in \mathcal{Z}} \sum_{k=0}^{23}
  \beta^{H}_{z,k}\, \mathbf{1}\{h(t)=k\}\, HDD_{z,d(t)} \\
  & + \alpha_{h(t),\,m(t)} + \gamma_{h(t),\,\tau(t)} + \delta_{y(t)} + \varepsilon_t .
\end{split}
\label{eq:zone}
\end{equation}

\subsection*{Indices}
\begin{itemize}
  \item $t$: an hourly observation (local time), 1 June 2021 to 31 December 2025,
        excluding 28--29 April 2025 (blackout); 40{,}148 hours.
  \item $h(t) \in \{0,\dots,23\}$: hour of the day of observation $t$.
  \item $d(t)$: calendar day of observation $t$.
  \item $m(t) \in \{1,\dots,12\}$: month of the year of $t$.
  \item $\tau(t)$: day type of $t$: Monday, \dots, Sunday, or holiday (national
        holidays and Holy Thursday form their own type).
  \item $y(t) \in \{2021,\dots,2025\}$: year of $t$.
  \item $z \in \mathcal{Z} = \{\text{North}, \text{Interior}, \text{Mediterranean}, \text{South}\}$: climate zone.
  \item $k \in \{0,\dots,23\}$: hour label in the sums.
\end{itemize}

\subsection*{Variables}
\begin{itemize}
  \item $D_t$: peninsular electricity demand in hour $t$ (MW, hourly mean).
  \item $\mathbf{1}\{h(t)=k\}$: indicator equal to 1 if observation $t$ falls in
        hour $k$ of the day and 0 otherwise. For each observation exactly one of
        the 24 indicators is 1, so only the coefficients of that hour apply.
  \item $CDD_{p,d} = \max(T_{p,d} - 21^{\circ}\mathrm{C},\, 0)$ and
        $HDD_{p,d} = \max(14^{\circ}\mathrm{C} - T_{p,d},\, 0)$: daily cooling and
        heating degree-days of province $p$ on day $d$ ($^{\circ}$C$\cdot$day),
        where $T_{p,d}$ is daily mean temperature. Because weather is daily, all
        24 hours of a day share the same value.
  \item $CDD_{z,d}$ and $HDD_{z,d}$: zone degree-days, the weighted mean over the
        provinces $p$ in zone $z$,
        \[
          CDD_{z,d} = \frac{\sum_{p \in z} w_p\, CDD_{p,d}}{\sum_{p \in z} w_p},
          \qquad
          w_p = s_{c(p)} \times \frac{\bar{q}_p}{\sum_{p' \in c(p)} \bar{q}_{p'}},
        \]
        where $s_{c}$ is community $c$'s share of peninsular demand over the
        sample and $\bar{q}_p$ is mean small-consumer demand in province $p$.
        The weights $w_p$ sum to one across the 47 provinces.
\end{itemize}

\subsection*{Coefficients and fixed effects}
\begin{itemize}
  \item $\beta^{C}_{z,k}$: effect of one additional cooling degree-day in zone $z$
        on log national demand at hour $k$, holding the other zones' weather
        fixed; $100\,\beta^{C}_{z,k}$ is approximately the percent change.
        $\beta^{H}_{z,k}$ is the analogue for heating degree-days. There are
        $4 \times 24 \times 2 = 192$ temperature coefficients.
  \item $\alpha_{h,m}$: hour-of-day $\times$ month fixed effects (288 cells);
        absorb the normal daily load shape in each month.
  \item $\gamma_{h,\tau}$: hour-of-day $\times$ day-type fixed effects (192 cells);
        absorb the weekly cycle and the holiday load shape.
  \item $\delta_{y}$: year fixed effects; absorb slow changes in the level of
        demand (growth, self-consumption, the 2022 price crisis).
  \item $\varepsilon_t$: error term. Standard errors are clustered by ISO week
        (240 clusters).
\end{itemize}

\subsection*{Climate zones}
\begin{center}
\begin{tabular}{@{}l >{\raggedright\arraybackslash}p{8.5cm} r r@{}}
\toprule
Zone & Autonomous communities & Provinces & Demand share \\
\midrule
North         & Galicia, Asturias, Cantabria, Pa\'is Vasco & 9 & 17.5\% \\
Interior      & Castilla y Le\'on, La Rioja, Navarra, Arag\'on, Madrid,
                Castilla-La Mancha, Extremadura & 22 & 31.5\% \\
Mediterranean & Catalu\~na, Comunidad Valenciana, Murcia & 8 & 34.3\% \\
South         & Andaluc\'ia & 8 & 16.7\% \\
\bottomrule
\end{tabular}
\end{center}

\subsection*{Example}
For the observation 15 July 2023 at 16:00, $h(t)=16$, so
$\mathbf{1}\{h(t)=16\}=1$ and all other indicators are 0. Equation
\eqref{eq:zone} reduces to
\[
\ln D_t = \sum_{z} \beta^{C}_{z,16}\, CDD_{z,\text{15 Jul}}
        + \sum_{z} \beta^{H}_{z,16}\, HDD_{z,\text{15 Jul}}
        + \alpha_{16,\text{Jul}} + \gamma_{16,\text{Sat}} + \delta_{2023} + \varepsilon_t .
\]

\end{document}
